\documentclass[10pt,a4paper]{amsart}
\usepackage[margin=19mm]{geometry}
\usepackage{amsmath,amssymb,mathtools}
\usepackage[scr=rsfs,cal=euler]{mathalfa}
\usepackage{xfrac}
\usepackage[unicode,hidelinks,bookmarks=false]{hyperref}
\newcommand{\ie}{i.e.\ }
\newcommand{\cf}{cf.\ }
\newcommand{\resp}{respectively\ }
\newcommand{\Subsection}{\subsection}
\providecommand{\nobreakdash}{\nobreak\hbox{-}}
\setlength{\emergencystretch}{2em}
\multlinegap0pt
\usepackage{bm}
\usepackage{bbm}
\usepackage{dsfont}
\usepackage{xspace}
\usepackage{cancel}
\usepackage{mathtools}
\usepackage{amsthm}
\makeatletter
\@ifundefined{thm}{\newtheorem{thm}{Theorem}}{}
\makeatother
\title[Weighted simple games and the topology of simplicial complex]{Weighted simple games and the topology of simplicial complexes}
\author{Anastasia Brooks, Franjo Sarcevic, Ismar Volic}
\date{Source: arXiv:2210.09771v3 (CC BY 4.0)}
\begin{document}
\maketitle

\noindent\emph{Source and licence.} Weighted simple games and the topology of simplicial complexes. Version arXiv:2210.09771v3 (CC BY 4.0), distributed under the Creative Commons Attribution 4.0 International licence, as recorded in the arXiv licence metadata for this exact version and shown on the abstract page https://arxiv.org/abs/2210.09771v3. Full rights notice: licenses/arxiv-2210.09771-CC-BY-4.0.txt.\par\smallskip

\noindent\textit{Editorial scope.} Counted unit: the Thm whose source label is \texttt{None} in the article (source0.tex lines 1144--1147, source label \texttt{None}), delivered together with the article's own complete original proof of it (source0.tex lines 1149--1183, one \texttt{proof} environment of 35 lines). No mathematics is written, repaired or re-derived: statement and proof are the article's own text line for line. Required context retained uncounted: none. Every definition, display and label of those lines is retained unchanged. Omitted from this package and not redistributed: 1--1143, 1148--1148, 1184--1880. Nothing inside a retained excerpt is omitted or altered: every retained body is delivered exactly as the article's lines are written, so no omission receipt of any kind accompanies this package. Citations: the retained excerpts carry no citation command at all, so no bibliography entry of the article is transcribed and no reference is redistributed. Reference adaptation: none was needed and none was made; every reference inside the delivered unit resolves inside it, so no unresolved reference or citation is produced. Printed slips kept literal: no printed word, symbol, hypothesis or number of the counted statement or of its proof was altered. Layout and class: the article's own class is \texttt{amsart} with packages including geometry, amssymb, amsmath, amscd, amsthm, color, epsfig, url, tikz, graphicx, xy; the wrapper is the lane's pinned \texttt{amsart} editorial wrapper at 10pt on A4, which restates the theorem-like environments the retained text needs and adds the article's own macro definitions that the retained text needs (none), with extra wrapper packages (bm, bbm, dsfont, xspace, cancel, mathtools). The delivered numbering is the unit's own. Assets: no retained span contains a figure environment, a graphics-inclusion command, a PDF-page inclusion, a table or a TikZ picture; no image file is copied into this package, the package declares no asset dependency and no image file is redistributed with it. Licence: version arXiv:2210.09771v3 of this article is distributed under the Creative Commons Attribution 4.0 International licence, as recorded in the arXiv licence metadata for this exact version and shown on the abstract page https://arxiv.org/abs/2210.09771v3. A full rights notice is delivered at licenses/arxiv-2210.09771-CC-BY-4.0.txt. Characters: every retained excerpt is pure ASCII, so the delivered engine, XeLaTeX with its default Latin Modern text font, reports no missing character.
\begin{thm}
Let $\{f_1,...,f_m\}$ be the set of all maximal simplices that contain $v_i$. For a maximal simplex $f_k$, let $P_{j}^{f_k, v_i}$ be the set of all $j$-dimensional faces of $f_k$ such that for an arbitrary $\varphi \in P_{j}^{f_k, v_i}$ no $(j-1)$-dimensional face of $\varphi$ in $\text{st}(v_i)$ reaches the quota, and let $n_{f_k}$ denote the number of vertices in $f_k$. Then the Shapley-Shubik index of the voter $v_i \in V$ in a weighted game $V(q;w_1,...,w_n)$ is given by
$$SS(v_i)=\frac{\sum_{k=1}^{m} \sum_{j=0}^{n_{f_k}-1} |\{s\in P_{j}^{f_k, v_i} \cap \text{st}(v_i): W_s \geq q\}|\cdot j! \cdot (n_{f_k}-(j+1))!}{\sum_{j=0}^d |\{s\in M_j: W_s \geq q\}|\cdot (j+1)!}.$$
\end{thm}

\begin{proof}
We can represent the maximal winning coalitions as the set of all maximal simplices in $K$ that contain vertices whose weights sum up to $q$ or greater. So the number of maximal winning coalitions is given by
$$\sum_{j=0}^d |\{s\in M_j: W_s \geq q\}|.$$
%where $d=\dim K$, $M_j$ is the set of all maximal simplices with $j+1$ vertices, and $W_s$ is the sum of the weights of the vertices of $s$.

The denominator for the Shapley-Shubik index, being the number of all maximal sequential coalitions, is then given by
$$\sum_{j=0}^d |\{s\in M_j: W_s \geq q\}|\cdot (j+1)!.$$
%\ifn{As usual, this should be presented as a proposition/theorem and a proof. Probably theorem since it has a substantial proof.}

The numerator for the Shapley-Shubik index of a voter $v_i$ in $V$ is the number of times $v_i$ is pivotal. Since the Shapley-Shubik index only considers maximal winning coalitions, in order to compute the index we can break down our simplicial complex $K$ into its maximal simplices and work through each maximal simplex separately. 

Note that for each maximal simplex $f=\{v_{j_1},...,v_{j_{n_f}}\}$, where $n_f$ is the number of vertices in $f$, a voter $v_i \in f$ is pivotal if the following is satisfied:
\begin{enumerate}
\item[(i)] $i=j_k$ for a $k\in \{1,...,n_f\};$
\item[(ii)]  the total weight of the face $\{v_{j_1},...,v_{j_k-1}\}$ does not reach the quota, i.e.
$$w_{j_1}+\cdots+w_{j_k-1}<q;$$
\item[(iii)]  the total weight of the face $\{v_{j_1},...,v_{j_k}\}$ reaches the quota, i.e.
$$w_{j_1}+\cdots+w_{j_k}\geq q.$$
\end{enumerate}

Hence we must only count the faces of $f$ containing $v_i$ that have weight greater or equal to $q$, but would have weight less than $q$ if we removed $v_i$. So the number of times $v_i$ is pivotal in a maximal sequential coalition $f$ is given by
$$\sum_{j=0}^{n_f-1} |\{s\in P_{j}^{f, v_i} \cap \text{st}(v_i): W_s \geq q\}|\cdot j! \cdot (n_f-(j+1))!,$$
where $P_{j}^{f, v_i}$ is the set of all $j$-dimensional faces of the maximal simplex $f$ such that for an arbitrary $\varphi \in P_{j}^{f, v_i}$, no $(j-1)$-dimensional face of $\varphi$ in $\text{st}(v_i)$ reaches the quota.
We multiply by $j!$ and $(n_f-(j+1))!$ to count the number of orders in which voters join a coalition before and after $v_i$.

To get the total number of times $v_i$ is pivotal in $V$, we evaluate the above expression for each maximal simplex that contains $v_i$ and take the sum of the resulting value. Thus, if $\{f_1,...,f_m\}$ is the set of all maximal simplices that contain $v_i$, the numerator for the Shapley-Shubik index of $v_i$ equals
$$\sum_{k=1}^{m} \sum_{j=0}^{n_{f_k}-1} |\{s\in P_{j}^{f_k, v_i} \cap \text{st}(v_i): W_s \geq q\}|\cdot j! \cdot (n_{f_k}-(j+1))!$$
%where $P_{j}^{f_k, v_i}$ is the set of all $j$-dimensional faces of the maximal simplex $f_k$ such that for an arbitrary $\varphi \in P_{j}^{f_k, v_i}$, no $(j-1)$-dimensional face of $\varphi$ in %\ifn{Does ``out of'' mean ``in''? I would just say ``in'' if it does.} 
%$\text{st}(v_i)$ reaches the quota, and $n_{f_k}$ is the number of vertices in $f_k$.
as desired.
%
%Therefore, the Shapley-Shubik index of the voter $v_i$ is given by
%$$SS(v_i)=\frac{\sum_{k=1}^{m} \sum_{j=0}^{n_{f_k}-1} |\{s\in P_{j}^{f_k, v_i} \cap \text{st}(v_i): W_s \geq q\}|\cdot j! \cdot (n_{f_k}-(j+1))!}{\sum_{j=0}^d |\{s\in M_j: W_s \geq q\}|\cdot (j+1)!}.$$
%
\end{proof}

\end{document}
